Fix one treatment-noise transform \(M\) represented by a law in
\(\mathcal P_{\mathrm{NG},n}\), and abbreviate
\[
 \psi=\psi_{\eta},\qquad
 R_{\mathrm{fac}}=R_1+1,\qquad
 R_{\mathrm{out}}=R_1+2.
\]
The certified intervals returned at error one contain their represented
values and have width at most one.  Consequently the integers in
\(\text{def:contour-bank-handle}\) satisfy
\[
 U_{\psi}\ge\max\{1,\psi\},
 \qquad U_R\ge\max\{1,R_1\}.                                      \tag{1}
\]
All calls used to obtain these integers terminate by the supplied modulus
contracts in \(\text{def:certified-contour-arithmetic-substrate}\).

We first bound the number of relevant zeros.  Membership in
\(\mathcal P_{\mathrm{NG},n}\) gives the exact Luxemburg inequality
\(E\exp(\eta^2/\psi^2)\le2\) and \(E\eta=0\).  The symmetrization and
Young-inequality calculation used in \(\text{lem:zero-localization}\) gives,
for every \(z\in\mathbb C\),
\[
 |M(z)|\le \exp(4\psi^2|z|^2).                                    \tag{2}
\]
Let \(N\) be the number of zeros of \(M\) strictly inside
\(\{|z|<R_{\mathrm{fac}}\}\), counted with multiplicity.  Jensen's formula,
applied at \(R_{\mathrm{out}}\) (with its usual radial limiting
interpretation if a zero lies on that circle), and \(M(0)=1\) imply
\[
 N\log\!\left(\frac{R_{\mathrm{out}}}{R_{\mathrm{fac}}}\right)
 \le 4\psi^2R_{\mathrm{out}}^2.                                   \tag{3}
\]
For \(x>0\), \(\log(1+x)\ge x/(1+x)\).  Taking
\(x=R_{\mathrm{fac}}^{-1}\), then using (1), gives
\[
 N\le4\psi^2R_{\mathrm{out}}^3
 \le4U_{\psi}^2(U_R+2)^3=N_{\mathrm{cert}}.                       \tag{4}
\]

Now put \(m=m_{\star}=N_{\mathrm{cert}}+3\),
\(d=2^{-m}\), and \(h=d/3\).  The grid in
\(\text{def:contour-bank-handle}\) consists of
\[
 \rho_q=R_0+\frac14+q2^{-m},
 \qquad 0\le q\le2^{m-1}.
\]
It has \(J=2^{m-1}+1>N_{\mathrm{cert}}\) points, spacing \(d\), and
\[
 R_0<\rho_q\le R_0+\frac34<R_1.                                  \tag{5}
\]
The modulus of a fixed zero can be at distance less than \(h\) from at
most one grid point, because \(2h<d\).  Every zero that can exclude a grid
point lies inside \(R_{\mathrm{fac}}\), since \(\rho_q+h<R_1+1\).  By (4) and the pigeonhole
principle there is an index \(j\) such that
\[
 \bigl||a|-\rho_j\bigr|\ge h                                    \tag{6}
\]
for every zero \(a\) strictly inside \(R_{\mathrm{fac}}\).  By
\(\text{lem:zero-localization}\), \(M\) has a zero with modulus at most
\(R_0\).  This zero is strictly inside \(C_j\) by (5), while (6) says that
no zero lies on \(C_j\).  The argument principle therefore gives
\[
 N_{C_j}(M)\ge1.                                                   \tag{7}
\]

It remains to verify the displayed dyadic modulus.  List all zeros strictly
inside \(R_{\mathrm{fac}}\), with multiplicity, as
\(a_1,\ldots,a_N\), and form the disk Blaschke product
\[
 Q(z)=\prod_{\nu=1}^{N}
 \frac{R_{\mathrm{fac}}(z-a_\nu)}
      {R_{\mathrm{fac}}^2-\overline{a_\nu}z}.
\]
After removable singularities are filled,
\(f=M/Q\) is holomorphic and zero-free on the open disk of radius
\(R_{\mathrm{fac}}\).  Each factor has modulus one on the boundary.
Thus (2) and the maximum-modulus principle give
\[
 |f(z)|\le \exp(4\psi^2R_{\mathrm{fac}}^2),
 \qquad |z|<R_{\mathrm{fac}}.                                    \tag{8}
\]
Moreover \(|Q(0)|\le1\) and \(M(0)=1\), so \(|f(0)|\ge1\).

Set \(A=4\psi^2R_{\mathrm{fac}}^2\).  The function
\(A-\log|f|\) is nonnegative and harmonic on the open disk, and its value
at zero is at most \(A\).  Harnack's inequality therefore yields, for
\(|z|\le R_1\),
\[
 \log|f(z)|
 \ge-A\left{
 \frac{R_{\mathrm{fac}}+|z|}{R_{\mathrm{fac}}-|z|}-1
 \right}
 \ge-8\psi^2R_1(R_1+1)^2.                                       \tag{9}
\]
With
\[
 E_{\star}=8U_RU_{\psi}^2(U_R+1)^2,
 \qquad D_{\star}=3(2U_R+1),                                    \tag{10}
\]
equations (1) and (9) imply \(|f(z)|\ge e^{-E_{\star}}\).  Since
\(\log2>1/2\),
\[
 |f(z)|\ge2^{-2E_{\star}},
 \qquad |z|\le R_1.                                              \tag{11}
\]

For \(z\in C_j\), (6) and \(|a_\nu|<R_{\mathrm{fac}}\) give, factor by
factor,
\[
 \left|
 \frac{R_{\mathrm{fac}}(z-a_\nu)}
      {R_{\mathrm{fac}}^2-\overline{a_\nu}z}
 \right|
 \ge \frac{|z-a_\nu|}{R_{\mathrm{fac}}+|z|}
 \ge \frac{h}{2R_1+1}
 \ge \frac{2^{-m}}{D_{\star}}.                                 \tag{12}
\]
Because \(D_{\star}\) is a positive integer,
\(D_{\star}^{-1}\ge2^{-D_{\star}}\).  The base in (12) is at most one;
hence (4) implies
\[
 |Q(z)|
 \ge2^{-N_{\mathrm{cert}}(m+D_{\star})}.                         \tag{13}
\]
Multiplying (11) and (13) gives
\[
 \inf_{z\in C_j}|M(z)|
 \ge2^{-\{2E_{\star}+N_{\mathrm{cert}}(m+D_{\star})\}}
 =2^{-u_{\star}}=a_{\star}>0.                                  \tag{14}
\]
The quantities \(E_{\star}\), \(D_{\star}\), and \(u_{\star}\) in
(10)--(14) are exactly the integers denoted by
\(e_{\star}\), \(d_{\star}^{\mathrm{den}}\), and \(u_{\star}\) in
\(\text{def:contour-bank-handle}\).

Finally, the construction makes only the two displayed modulus calls,
certified additions with rational point records, and finite primitive
recursions over explicit integer bounds.  Soundness follows from
\(\text{def:certified-contour-arithmetic-substrate}\), and no operation
queries \(M\), \(\eta\), or their law.  The input record is fixed before
\(n\), while (2)--(14) use only the fixed class constants.  The same finite
bank and the same positive dyadic certificate therefore work uniformly in
\(n\).